\section{Discussion}

This paper provides evidence that microstructure-based measures of trading pressure contain economically meaningful information about post-earnings return dynamics. In particular, we find that persistent and directional order flow prior to earnings announcements is associated with attenuated post-announcement drift, consistent with partial information incorporation before the event.

\subsection{Interpretation and Economic Mechanism}

The central interpretation of our findings is that not all price adjustment occurs at the time of the earnings announcement. Instead, information appears to be incorporated gradually through trading activity in the pre-announcement period. When order flow reflects informed trading—characterized by persistent and directional imbalances—prices adjust partially before the public release of information, leaving less room for subsequent drift.

This interpretation is consistent with canonical models of informed trading, such as those of \citet{Kyle1985} and \citet{Hasbrouck1991}, in which private information is revealed through order flow and gradually impounded into prices. In this framework, the magnitude of post-announcement drift depends not only on the content of the earnings news, but also on the extent to which that information has already been incorporated into prices.

Our results also highlight that different dimensions of trading activity play distinct roles. While both the information factor and liquidity-related components are associated with returns, the stronger and more consistent effect of the information factor suggests that persistent directional trading is the primary channel through which pre-announcement activity affects post-event price dynamics.

\subsection{Relation to the PEAD Literature}

The findings contribute to the post-earnings-announcement drift (PEAD) literature by introducing a microstructure-based perspective. Traditional approaches measure news using accounting-based signals such as standardized unexpected earnings, implicitly treating the announcement as the primary moment of information revelation. In contrast, our results suggest that the timing of information incorporation varies across events and can be partially observed through trading behavior prior to the announcement.

By showing that microstructure-based measures retain predictive power even after controlling for SUE, we provide evidence that trading activity captures an additional dimension of information beyond the magnitude of earnings surprises. This suggests that PEAD may, in part, reflect heterogeneity in the speed of information diffusion rather than solely delayed reaction to public news.

\subsection{Limitations}

Several limitations should be noted. First, the decomposition of trading pressure into information-driven and liquidity-driven components is based on a reduced-form approach. While the PCA-based factors and pressure score are economically interpretable, they do not provide a structural identification of informed trading. As a result, alternative explanations—such as temporary price pressure or residual liquidity effects—cannot be fully ruled out.

Second, the explanatory power of the models is modest, with low $R^2$ values typical of cross-sectional return regressions. This reflects the inherently noisy nature of financial return data and suggests that microstructure-based signals capture only one component of a broader set of drivers of post-earnings returns.

Third, measurement error in high-frequency data may affect the construction of microstructure features. Although we apply standard cleaning and aggregation procedures, TAQ data are subject to classification errors and microstructure noise that may attenuate estimated relationships.

\subsection{Implications and Extensions}

Despite these limitations, the results have several implications. First, they suggest that incorporating high-frequency trading information can enhance our understanding of return predictability around earnings announcements. Second, they highlight the importance of distinguishing between the magnitude of information (as measured by earnings surprises) and the timing of information incorporation (as reflected in trading activity).

Several extensions could further strengthen the analysis. One direction is to examine intraday dynamics more directly, including the evolution of order flow and price impact within the pre-announcement window. Another is to incorporate richer measures of liquidity and market depth, which may help disentangle information-driven trading from liquidity shocks. Finally, applying the framework to other information events, such as macroeconomic announcements or corporate disclosures, could test the generality of the findings.

\subsection{Conclusion of Findings}

Overall, the evidence suggests that pre-announcement trading activity contains valuable information about the subsequent evolution of prices. By combining microstructure features with traditional event-study methods, this paper provides a more complete picture of how information is incorporated into asset prices around earnings announcements.

